
\section{Early Attempts at Solving Sokoban}

In the late 1990s and 2000s, the focus of Sokoban research shifted from the difficulty of the problem to its solution. Early attempts at solving Sokoban relied on variations of traditional algorithmic search such as A*. Junghanns et al.\ presented numerous research projects on Sokoban in the late 1990s. In one of these works, they attempted to solve all 90 puzzles of the XSokoban test suite using IDA* \cite{myers_xsokoban, junghanns1997sokoban, junghanns1998sokoban}. The researchers demonstrated a modified IDA* search method named ``Rolling Stone'' that incorporates a Deadlock Table to identify dead-end configurations earlier in the heuristic search-tree expansion in order to prune the search. The Deadlock Table consisted of 22 million entries of frequently occurring board patterns that lead to a dead end. Rolling Stone also introduced a concept called ``Macro Moves'' that collapses a sequence of forced moves into single steps for more efficient memory use during search while preserving solution optimality.

The Rolling Stone algorithm was able to find solutions to 16 of the 90 problems with a search tree shallower than the solution itself. This was made possible by Macro Moves. However, this method comes with significant limitations, as it requires a 22-million-row deadlock pattern database, a piece of hard-coded domain knowledge. The paper also recognizes that Rolling Stone spends 90\% of its execution time updating the lower bounds, which means most of the processing time is spent on optimizing a found solution \cite{junghanns1998sokoban}.

A similar method was tried by Pereira et al.\ \cite{pereira2013finding}, where the researchers employed a pattern database, a pre-computed lookup table that stores the shortest distances for simplified local problems, to be used as a heuristic in the IDA* search. The authors introduced a concept called instance decomposition. Because it is not possible to know which stone is placed on which goal position before the puzzle is solved, instance decomposition assigns each stone to a specific goal before running the search. This constraint acts as an intermediate goal and converts the multi-goal problem into a single-goal problem. During the search, the Pattern Database provides an admissible heuristic to be used for pruning the search tree. This motivates the search for better heuristics and representations.